\documentclass{beamer}
\mode<presentation> {\usetheme{default}}
%\setbeamersize{text margin left=10mm,text margin right=10mm} 

\setbeamertemplate{footline}[frame number]


\usepackage{amssymb}
\usepackage{multirow}
\usepackage{amsthm}
\usepackage{amsmath}
\usepackage{booktabs}
\usepackage{bm}
\usepackage{needspace}
\usepackage{comment}
\usepackage{tikz}
\usepackage{graphicx}
\usepackage{bbm, dsfont}

\usetikzlibrary{arrows,shapes,calc,positioning}

\theoremstyle{plain}
\newtheorem{thm}{Theorem}[section]
\newtheorem{cor}[thm]{Corollary}
\newtheorem{lem}[thm]{Lemma}
\newtheorem{prop}[thm]{Proposition}
%
\theoremstyle{definition}
\newtheorem{defn}{Definition}[section]
\newtheorem{remk}{Definition}[section]
\newtheorem{exam}{Example}[section]
\newtheorem{prob}{Exercise}[section]
\newtheorem{solu2}{Solution}[section]
\newenvironment{Shaded}{\begin{snugshade}}{\end{snugshade}}

\newcommand{\dif}{\operatorname{d}}
\newcommand{\var}{\operatorname{Var}}

%opening
\title{Linear Mixed Models}
\author{Himchan Jeong, FSA}

\date{}

\begin{document}
\maketitle

\section{Linear Mixed Models}

\subsection{Motivation}

\begin{frame}[allowframebreaks]
\frametitle{Motivation}

Recall the basic assumptions of the linear models where $y_1,\ldots,y_m$ are the observed response variable and $\mathbf{x}_1,\ldots, \mathbf{x}_m \in \mathbb{R}^{p \times 1}$ are the individual covariates  (here $Z \in \mathbb{R}^{a \times b}$ means that $Z$ is a $(a \times b)$ real-valued matrix):
\begin{enumerate}
	\item $E[y_i] = \mathbf{x}_i' \bm{\beta}$ for $\bm{\beta} = (\beta_1, \ldots, \beta_p) \in \mathbb{R}^{p\times 1}$,
	\item $ \{\mathbf{x}_1, \mathbf{x}_2, \ldots, \mathbf{x}_m \}$ are non-random,
	\item $Var(y_i) = \sigma^2$,
	\item $\{ y_1, y_2, \ldots, y_m  \}$ are independent,
	\item $\{ y_1, y_2, \ldots, y_m  \}$ are normally distributed.
\end{enumerate}
Such situations can be written in short as follows:
\begin{equation} \label{eq:linmod}
y_i = \mathbf{x}_i' \bm{\beta}  + \epsilon_i, \quad \epsilon_i \overset{i.i.d.}\sim \mathcal{N}(0, \sigma^2).
\end{equation}

Note that the first components of $ \{\mathbf{x}_1, \mathbf{x}_2, \ldots, \mathbf{x}_m \}$ are all 1 so that $\beta_1$ corresponds to the \emph{intercept} term in the linear models.
\framebreak

Now, let us assume an example of data with hierarchy so that we observe data points for each of the subjects $i \in \{1,\ldots, m\}$ for time $t=1,\ldots, n_i$. If we adopt the assumptions from the linear models, then we can write them as follows with the modified subscripts:

\begin{enumerate}
	\item $E[y_{i,t}] = \mathbf{x}_{i,t}' \bm{\beta}$ for $\beta \in \mathbb{R}^{p \times 1}$,
	\item $ \{\mathbf{x}_{1,1}, \mathbf{x}_{1,2}, \ldots, \mathbf{x}_{m,n_m} \}$ are non-random,
	\item $Var(y_{i,t}) = \sigma^2$,
	\item {\color{red} $\{ y_{1,1}, y_{1,2}, \ldots, y_{m,n_m}  \}$ are independent,}
	\item $\{  y_{1,1}, y_{1,2}, \ldots, y_{m,n_m}   \}$ are normally distributed.
\end{enumerate}
What if one of the above assumptions is not valid...? (Note that $\{ y_{i,1}, y_{i,2}, \ldots, y_{i,n_i}  \}$ are from the same subject $i$ over time)

\framebreak
In this regard, one can consider the following representation:
\begin{equation} \label{eq:mixedmod_intercept}
y_{i,t} = \mathbf{x}_{i,t}'\bm{\beta} + {\color{red}u_i}+ \epsilon_{i,t},
\end{equation}
where $\epsilon_{i,t} \overset{i.i.d.}\sim \mathcal{N}(0, \sigma^2)$  for $t=1,\ldots, n_i$ and $u_i \overset{i.i.d.}\sim \mathcal{N}(0, \phi_u)$ for $i=1,\ldots, m$.

\vspace{8pt}
\eqref{eq:mixedmod_intercept} is one of the most widely used example of the \textbf{linear mixed models}. In this case, $y_{i,t}$ and $y_{i,t'}$ are NOT independent!


\framebreak

Let us consider two special cases of \eqref{eq:mixedmod_intercept} here:
\begin{itemize}
	\item Complete pooling (when $\phi_u \rightarrow 0$): If so, $u_i \equiv 0$ so that it is reduced to \eqref{eq:linmod}.
	\item No pooling (when $\phi_u \rightarrow \infty$): If so, we can analyze the data only if for the data points from the same subject $i$.
\end{itemize}

Therefore, by allowing $\phi$ to be a positive constant, we can do multi-level modeling somewhere between the above two extreme cases; complete pooling (ignore group effects) and no pooling (all the groups are treated as different datasets).

\end{frame}


\begin{frame}{Why Random Effects?}

The following hypothetical example shows us that we might capture the
\color{purple} unobserved heterogeneity in risk \color{black} by
observing the residual after controlling for the effects of observed
covariates, which can be explained in terms of \color{purple} random
effects \color{black} for policyholders A and B.

\begin{table}[h!t!]
\centering
\resizebox{\textwidth}{!}{
\begin{tabular}{c|c|c|c|c|c|c|c}
\multirow{2}{*}{Year} & \multirow{2}{*}{Gender} & \multirow{2}{*}{Age} & \multirow{2}{*}{Vehicle Size} & \multicolumn{2}{c|}{Policyholder A} & \multicolumn{2}{c}{Policyholder B} \\ \cline{5-8} 
                      &                         &                      &                               & \# of Claim      & Amt of Claim     & \# of Claim      & Amt of Claim     \\ \hline
2015                  & M                       & 45                   & Medium                        & 0                & 0                & 1                & 500              \\ 
2016                  & M                       & 46                   & Medium                        & 0                & 0                & 2                & 4000             \\ 
2017                  & M                       & 47                   & Large                         & 1                & 200              & 1                & 8000             \\ 
\end{tabular}}
\end{table}
Note that random effects model has a good connection with  \color{purple} credibility theory  \color{black} as well.
\end{frame}



\subsection{Inference}

\begin{frame}[allowframebreaks]
\frametitle{Model assumptions and notation}

Now, we are ready to introduce a general form of the linear mixed models given as follows:

\begin{equation} \label{eq:mixedmodgen}
\begin{aligned}
&\bm{y} = \bm{X\beta + Zu + \epsilon} \\
&\bm{u} \sim \mathcal{N}(\bm{0, D}), \ \bm{\epsilon} \sim \mathcal{N}(\bm{0, \Sigma}), \ \bm{u} \perp \bm{\epsilon},
\end{aligned}
\end{equation}
where $\bm{Z}$ is a pre-determined model matrix, $\bm{D}$ is the covariance matrix of the random effects $\bm{u}$, and $\bm{\Sigma}$ is the covariance matrix of the white noise $\bm{\epsilon}$. 

Note that the marginal/conditional mean and variance of $\mathbf{y}$ are given as follows:
\begin{itemize}
	\item $E[\bm{y}|\bm{u}] = \bm{X\beta + Zu}$, $Var(\bm{y}|\bm{u}) = \bm{\Sigma}$,
	\item $E[\bm{y}] = \bm{X\beta}$, $Var(\bm{y}) = \bm{ZDZ' + \Sigma:= V}$ ({\color{red} why?}).
\end{itemize}

\framebreak

\begin{exam}[Random intercepts]
Recall the motivational example in \eqref{eq:mixedmod_intercept}. By letting $\bm{u}=(u_1, u_2, \ldots, u_m)'$, one can check that, it is a special case of \eqref{eq:mixedmodgen} where $\bm{\Sigma} = \sigma^2 \cdot \bm{I}_{N} \in \mathbb{R}^{N \times N}$ ($N = n_1 + \ldots+ n_m$), $\bm{D}= \phi_u \cdot \bm{I}_{{m}}\in \mathbb{R}^{{m} \times {m}}$, and
$$
\bm{Z}=\begin{bmatrix}
\bm{1}_{n_1} & \bm{0}_{n_1} & \cdots & \bm{0}_{n_1} \\
\bm{0}_{n_2} & \bm{1}_{n_2} &  &  \\
\vdots &  & \ddots &  \\
\bm{0}_{n_m} &  &  & \bm{1}_{n_m} \\
\end{bmatrix}\in \mathbb{R}^{N \times m}.
$$
Here $\bm{1}_{k}$ (resp. $\bm{0}_{k}$) means a vector of length $k$ with all components being 1 (resp. 0).
\end{exam}

\framebreak


More specifically, suppose we have
$$
y_{i,t} = \mathbf{x}_{i,t}'\bm{\beta} + {\color{red}u_i}+ \epsilon_{i,t},
$$
where $\epsilon_{i,t} \overset{i.i.d.}\sim \mathcal{N}(0, \sigma^2)$  for $t=1,2, 3$ and $u_i \overset{i.i.d.}\sim \mathcal{N}(0, \phi_u)$ for $i=1,2, 3$. In this case, we have $\bm{\Sigma} = \sigma^2 \cdot \bm{I}_{\bm{T}} = \sigma^2 \cdot \bm{I}_{9} \in \mathbb{R}^{9 \times 9}$, $\bm{D}= \phi_u \cdot \bm{I}_{{3}}\in \mathbb{R}^{{3} \times {3}}$, and
$$
\bm{Z}'=\begin{bmatrix}
1 & 1 & 1 & 0 & 0 & 0 & 0 & 0 & 0 \\
0 & 0 & 0 & 1 & 1 & 1 & 0 & 0 & 0 \\
0 & 0 & 0 & 0 & 0 & 0 & 1 & 1 & 1 \\
\end{bmatrix}\in \mathbb{R}^{3 \times 9}.
$$

\framebreak

\begin{exam}[Random slopes]
Suppose there are three clusters ($m=3$) where each of the clusters have 1000 data points ($n_1=n_2=n_3=1000$) with a covariate ${x}_{i,t} \in \mathbb{R}^{1 \times 1}$  and the corresponding response variable $y_{i,t}$. 

If we consider a situation where the {\color{red} impacts of the covariate are different among the clusters}, then it can be modeled as follows:
$$
y_{i,t} = \beta_1 + \beta_2 {x}_{i,t} + {\color{red}v_i {x}_{i,t}}+ \epsilon_{i,t}, \ i=1,2,3, \ t=1,\ldots, 1000,
$$
where $\epsilon_{i,t} \overset{i.i.d.}\sim \mathcal{N}(0, \sigma^2)$  and $v_i \overset{i.i.d.}\sim \mathcal{N}(0, \phi_v)$. In this case, we have $\bm{\Sigma} =  \sigma^2 \cdot \bm{I}_{3000} \in \mathbb{R}^{3000 \times 3000}$, 
$\bm{D}= \phi_v \cdot \bm{I}_3 \in \mathbb{R}^{{3} \times {3}}$, and %(Let $\bm{u}=(v_1, v_2, v_3)'$ here.) 
$$
{\small
\bm{Z}'=\begin{bmatrix}
 {x}_{1,1} & \cdots &  {x}_{1,1000} & 0 & \cdots & 0 & 0 & \cdots & 0 \\
0 & \cdots & 0 & {x}_{2,1} & \cdots &  {x}_{2,1000} & 0 & \cdots & 0 \\
0 & \cdots & 0 & 0 & \cdots & 0 & {x}_{3,1} & \cdots &  {x}_{3,1000}\\
\end{bmatrix}\in \mathbb{R}^{3 \times 3000}. }
$$
\end{exam}

\pagebreak

\begin{exam}[Random intercepts and slopes]
Suppose there are three clusters ($m=3$) where each of the clusters have 1000 data points ($n_1=n_2=n_3=1000$) with a covariate ${x}_{i,t} \in \mathbb{R}^{1 \times 1}$  and the corresponding response variable $y_{i,t}$. 

If we consider a situation where the {\color{red} impacts of the covariate are different among the clusters on top of the different intercepts}, then it can be modeled as follows:
$$
y_{i,t} = \beta_1 + \beta_2 {x}_{i,t} + {\color{red}u_i + v_i {x}_{i,t}}+ \epsilon_{i,t}, \ i=1,2,3, \ t=1,\ldots, 1000,
$$
where $\epsilon_{i,t} \overset{i.i.d.}\sim \mathcal{N}(0, \sigma^2)$  for $t=1,\ldots, 1000$, $u_i \overset{i.i.d.}\sim \mathcal{N}(0, \phi_u)$, $v_i \overset{i.i.d.}\sim \mathcal{N}(0, \phi_v)$ for $i=1,2,3$, and $u_i \perp v_i$.

\vspace{8pt}
What would be $\bm{\Sigma}$, $\bm{D}$, and $\bm{Z}$ in this case? (Let $\bm{u}=(u_1, u_2, u_3, v_1,v_2, v_3)'$ here.)

\end{exam}


\end{frame}


\begin{frame}[allowframebreaks]
\frametitle{Estimation and prediction}
\begin{theorem}[Estimation of $\bm{\beta}$ with known $\bm{V}$]
If $\bm{V = ZDZ' + \Sigma}$ is known and positive definite, and $X$ is of full rank, then the generalized least squares (GLS) estimator of $\bm{\beta}$ is given as follows:
$$
\hat{\bm{\beta}}^{GLS}=\bm{(X'V^{-1}X)^{-1} X'V^{-1}y}.
$$
\end{theorem}

\begin{proof}
See \url{https://www.statlect.com/fundamentals-of-statistics/generalized-least-squares}.
\end{proof}

\framebreak
For simplicity, let us define $\bm{C = (X'V^{-1}X)^{-1} X'V^{-1}}$. Then it is easy to check that
$$
\hat{\bm{\beta}}^{GLS} = \bm{C X\beta}  + \bm{C Zu} +  \bm{C \epsilon} = \bm{\beta}  + \bm{C Zu} +  \bm{C \epsilon},
$$
so that $\mathbb{E}[\hat{\bm{\beta}}^{GLS}]=\bm{\beta} + \bm{C Z \mathbb{E}[u]} + \bm{C  \mathbb{E}[\epsilon]=\bm\beta} $ (unbiased), and
$$
\begin{aligned}
\mathbb{V}ar[\hat{\bm{\beta}}^{GLS}] &= \bm{C Z \mathbb{V}ar[u] Z' C'} + \bm{C  \mathbb{V}ar[\epsilon] C'} \\
 &=  \bm{C Z D Z' C'} + \bm{C  \Sigma C'}  = \bm{C  V C'}\\
& = \bm{(X'V^{-1}X)^{-1} X'V^{-1} \cdot V  \cdot V^{-1} X(X'V^{-1}X)^{-1} }\\
&=\bm{(X'V^{-1}X)^{-1}}.
\end{aligned}
$$

What about $\hat{\bm{\beta}}^{OLS}=\bm{(X'X)^{-1}X'y}$?

\framebreak

\begin{theorem}[Prediction of $\bm{u}$ with known $\bm{V}$ and $\bm{\beta}$]
The best linear predictor (BLP) of $\bm{u}$ is given as follows:
$$
\text{BLP}(\bm{u}):=\hat{\bm{u}}=\bm{D Z'V^{-1}(y-X\beta)}.
$$
\end{theorem}

\begin{proof}
See page 197 of the textbook (Chapter 8 of PMAS1).
\end{proof}
Usually, $\bm{\beta}$ is unknown so it is replaced with $\hat{\bm{\beta}}$ in practice in that case. Note also that $\hat{\bm{u}}$ is also an unbiased estimator of $\bm{u}$, namely, $E[\hat{\bm{u}}]=E[\bm{u}]=0$.

\framebreak



By further (innocuously) assuming that $\bm{u}_i$ are i.i.d. (see the notation of (8.22) in the textbook), one can write out the BLP of $\bm{u}_i$ as follows:
$$
\hat{\bm{u}}_i=G Z_i'V_i^{-1}(\bm{y}_i-X_i\hat{\bm{\beta}}),
$$
where the defintion of $G, \bm{u}_i, Z_i, V_i, \bm{y}_i, X_i$ can be found from (8.20)--(8.24) of the textbook.

\vspace{8pt}

Suppose that for the subject $i$, we observed a new covariate vector $\bm{x}_{i,n_i+1}$. In this case, the predicted value of $\bm{y}_{i,n_i+1}$ is given as

$$
\hat{\bm{y}}_{i,n_i+1}= \bm{x}'_{i,n_i+1}\hat{\bm{\beta}} + \bm{z}'_{i,n_i+1}\hat{\bm{u}}_i.
$$

\framebreak

Let us consider the setting in \eqref{eq:mixedmod_intercept} where $\mathbf{x}_{i,t} = 1 \in \mathbb{R}^{1 \times 1}$ and $n_i=n$ so that
$$
y_{i,t} = \mu + {\color{red}u_i}+ \epsilon_{i,t},
$$
where $\epsilon_{i,t} \overset{i.i.d.}\sim \mathcal{N}(0, \sigma^2)$  for $t=1,\ldots, n$ and $u_i \overset{i.i.d.}\sim \mathcal{N}(0, \phi_u)$ for $i=1,\ldots, m$. In this case, it is easy to figure out that
$$
\begin{aligned}
G &= \phi_u \in \mathbb{R}^{1 \times 1}, Z_i = \bm{1}_{n} \in \mathbb{R}^{n \times 1}, {\Sigma}_i = \sigma^2 \cdot I_n \in \mathbb{R}^{n \times n}, \\
V_i &= Z_i G Z_i' + {\Sigma}_i = \phi_u \bm{1}_{n} \bm{1}_{n}' + \sigma^2 \cdot I_n \in \mathbb{R}^{n \times n} = \phi_u \cdot J_n + \sigma^2 \cdot I_n,
\end{aligned}
$$
where $J_n \in \mathbb{R}^{n \times n}$ is a matrix with all components being 1.

Question: How do we find $V_i^{-1}$ in this case?

Answer: Via \href{https://en.wikipedia.org/wiki/Sherman-Morrison_formula}{ `Sherman–Morrison formula'}.
$$
V_i^{-1}= \sigma^{-2} \cdot I_n - \frac{\sigma^{-4}\phi_u J_n}{1 + \sigma^{-2}\phi_u  n},
$$
because $ \bm{1}_{n} ' \bm{1}_{n} =n$.

\framebreak
Since $ \bm{1}_{n} ' J_n =n \bm{1}_n'$, $\mathbf{x}_{i,t}\bm{\beta}=\mu$, and $\bm{1}'_n \bm{y}_i = \sum_{t=1}^n y_{i,t}= n \overline{y}_i$, we have
$$
\begin{aligned}
\hat{\bm{y}}_{i,n_i+1} &= \bm{x}'_{i,n_i+1}\hat{\bm{\beta}} + \bm{z}'_{i,n_i+1}\hat{\bm{u}}_i =  \hat{\mu} + \hat{\bm{u}}_i \\
								 &= \hat{\mu} + G Z_i'V_i^{-1}(\bm{y}_i-X_i\hat{\bm{\beta}}) \\
								 &= \hat{\mu} + \phi_u  \bm{1}_{n} '\left( \sigma^{-2} \cdot I_n - \frac{\sigma^{-4}\phi_u J_n}{1 + \sigma^{-2}\phi_u  n} \right)(\bm{y}_i-\hat{\mu}\cdot \bm{1}_{n}) \\
								&= \hat{\mu} + \sigma^{-2}\phi_u   \left(\bm{1}_n' - \frac{\sigma^{-2}\phi_u n \bm{1}_n'}{1 + \sigma^{-2}\phi_u  n} \right)(\bm{y}_i-\hat{\mu}\cdot \bm{1}_{n}) \\
								&= \hat{\mu} + \sigma^{-2}\phi_u   \left(1 - \frac{\sigma^{-2}\phi_u n}{1 + \sigma^{-2}\phi_u  n} \right)(\bm{1}_n'\bm{y}_i-\hat{\mu}\cdot \bm{1}_n'\bm{1}_{n}) \\
								&= \hat{\mu} + \frac{\sigma^{-2}\phi_u n}{1 + \sigma^{-2}\phi_u  n} (\overline{y}_i-\hat{\mu}) \\
								&= \frac{\sigma^2/\phi_u}{\sigma^2/\phi_u+n}\hat{\mu} + \frac{n}{\sigma^2/\phi_u+n} \overline{y}_i = (1-Z) \hat{\mu} + Z \overline{y}_i.\\
\end{aligned}
$$
Later, it turns out that such a result is closely related with the so-called `credibility premium'.

\framebreak
Recall two special cases of \eqref{eq:mixedmod_intercept} here:
\begin{itemize}
	\item Complete pooling (when $\phi_u \rightarrow 0$): If so, $Z=\frac{\phi_u n}{\sigma^2+\phi_u n} \rightarrow  0$ so that 
$$
\hat{\bm{y}}_{i,n_i+1} = \frac{\sigma^2/\phi_u}{\sigma^2/\phi_u+n}\hat{\mu} + \frac{n}{\sigma^2/\phi_u+n} \overline{y}_i \rightarrow \hat{\mu}.
$$
	\item No pooling (when $\phi_u \rightarrow \infty$): If so, $Z=\frac{n}{\sigma^2/\phi_u+ n} \rightarrow  1$ so that
\end{itemize}
$$
\hat{\bm{y}}_{i,n_i+1} = \frac{\sigma^2/\phi_u}{\sigma^2/\phi_u+n}\hat{\mu} + \frac{n}{\sigma^2/\phi_u+n} \overline{y}_i \rightarrow  \overline{y}_i.
$$

Therefore, by allowing $\phi_u$ to be a positive constant, we can get a predicted value of $\hat{\bm{y}}_{i,n_i+1}$ in the spirit of credibility premium; which is a compromise of cluster experience, $\overline{y}_i$, and the (estimate of) manual premium, $\hat{\mu}$.

\framebreak

\begin{itemize}
	\item In practice, the covariance matrix $\bm{V(\theta)=ZD(\theta)Z' + \Sigma(\theta)}$ is often unknown so they need to be estimated from the data using the likelihood $L(\bm{\beta, \theta})$ as in (8.37) ($\bm{\theta}$ is the covariance parameters).
	\item If $\bm{\theta}$ is fixed, then the maximizer of (8.37) with respect to $\bm{\beta}$ is given as follows:
$$
\bm{\hat{\beta}(\theta) = (X'V(\theta)^{-1}X)^{-1} X'V(\theta)^{-1}y}.
$$
That being said, if one can obtain a reasonable estimate of $\bm{\theta}$, then the joint estimation of $\bm{\beta, \theta}$ is done.
\end{itemize}

\framebreak
\begin{itemize}
	\item Restricted maximum likelihood (REML) is a popular way to estimate $\bm{\theta}$ for the LMMs, where we use the following marginal log-likelihood of $\bm{\theta}$ by integrating out the fixed effects in $\bm{\beta}$:
$$
\begin{aligned}
\ell_r(\bm{\theta})&= \ln  \left( \int L(\bm{\beta, \theta})d \bm{\beta} \right) \\
					&= \ell_p(\bm{\theta}) - \frac{1}{2} \ln |\bm{X'V(\theta)^{-1}X}| + \text{constants},
\end{aligned}
$$
where $\ell_p(\bm{\theta})$ is defined in (8.39) of the textbook.
\item In \texttt{R}, one can choose whether to use REML approach (minimizing $\ell_r(\bm{\theta})$) or ML approach (minimizing $\ell_p(\bm{\theta})$) with the \texttt{lmer} function by specifying \texttt{REML=TRUE} (by default) or \texttt{REML=FALSE} (recall $\bm{y}'M_X \bm{y} / n$ vs. $\bm{y}'M_X \bm{y} / (n-p)$ in the linear models).
\end{itemize}


\end{frame}


%\begin{theorem}[Prediction of $\bm{y}$ with known $\bm{V}$ and updated $\bm{X}_{i,*}$]
%Suppose that for the subject $i$, we observed a new covariate vector $\bm{X}_{i,*}$. In this case, the predicted value of $\bm{y}_{i,*}$ is given as follows:
%$$
%\begin{aligned}
%\hat{\bm{y}}_{i,*}:&= \bm{X}_{i,*}\hat{\bm{\beta}} + \bm{Z}_i\hat{\bm{u}}_i \\
%			  &=  \bm{X}_{i,*}\hat{\bm{\beta}} + \bm{Z}_i\bm{G Z_i'V_i^{-1}(y_i-X_{i}\hat{\bm{\beta}})} \\
%			  &=  \bm{X}_{i,*}\hat{\bm{\beta}} + (\bm{Z_iG Z_i'+\Sigma_i})\bm{V_i^{-1}(y_i-X_{i}\hat{\bm{\beta}})} - \bm{\Sigma_i}\bm{V_i^{-1}(y_i-X_{i}\hat{\bm{\beta}} )}\\
%&= \bm{(\Sigma_i V_i^{-1}) \cdot X_{i,*}\hat{\bm{\beta}}} + (\bm{I_{T_i} - \Sigma_i V_i^{-1}}) \cdot \bm{y}_i
%\end{aligned}
%$$
%\end{theorem}



\begin{frame}[allowframebreaks]
\frametitle{Applications}

\begin{exam}[Workers’ Compensation Insurance Losses]
We want to predict the worker's compensation insurance losses as a function of year and risk-class where the payroll is given as an offset. Let us consider the following linear models first (no random effects):
\begin{itemize}
	\item Complete pooling: clustering of data in risk classes are not considered at all
$$
\log (\texttt{Loss}_{ij}) = \log (\texttt{Payroll}_{ij}) +\beta_1 + \beta_2 \texttt{Year}_{ij}+ \epsilon_{ij}.
$$
	\item No pooling: clustering of data in risk classes are considered as fixed effects
$$
\log (\texttt{Loss}_{ij}) = \log (\texttt{Payroll}_{ij}) +\beta_{1, {\color{red}i}} + \beta_2 \texttt{Year}_{ij}+ \epsilon_{ij}.
$$
\end{itemize}

\end{exam}

\framebreak

One can consider the following linear mixed models as well:
\begin{itemize}
	\item LMM with random intercepts: 
$$
\log (\texttt{Loss}_{ij}) = \log (\texttt{Payroll}_{ij}) +\beta_1 + {\color{red}u_i}+ \beta_2 \texttt{Year}_{ij}+ \epsilon_{ij}.
$$
	\item LMM with random intercepts and slopes: 
$$
\log (\texttt{Loss}_{ij}) = \log (\texttt{Payroll}_{ij}) +\beta_1 + {\color{red}u_i}+ \beta_2 \texttt{Year}_{ij}+{\color{red}v_i} \texttt{Year}_{ij}+ \epsilon_{ij}.
$$
\end{itemize}

\end{frame}


\end{document}